\section{Circuit Complexity --- Circuit Sat Reduction Valid}
%%% generated by the blueprint pipeline from TCSlib.Complexity.CircuitComplexity.CircuitSatReductionValid
\subsection*{Circuit Complexity --- Circuit Sat Reduction Valid: Introduction}

Declarations of the Lean module \texttt{TCSlib.Complexity.CircuitComplexity.CircuitSatReductionValid}.

\subsection*{Circuit Complexity --- Circuit Sat Reduction Valid: Results}

\begin{definition}[Dual gate kind]
\lean{BoolCircuit.GateKind.dual}
\label{BoolCircuit.GateKind.dual}
\leanok
\uses{BoolCircuit.GateKind}
The dual of a gate kind interchanges conjunction and disjunction and leaves negation
unchanged.
\end{definition}

\begin{definition}[Dual gate]
\lean{BoolCircuit.DAGGate.dual}
\label{BoolCircuit.DAGGate.dual}
\leanok
\uses{BoolCircuit.DAGGate, BoolCircuit.GateKind.dual}
The dual of a directed acyclic graph gate has the dual gate kind and the same list of
argument vertices as the original gate.
\end{definition}

\begin{theorem}[Disjunction and negated conjunction]
\lean{BoolCircuit.any_eq_not_all}
\label{BoolCircuit.any_eq_not_all}
\leanok
\difficulty{4}
Let $l$ be a list of natural numbers, and let $p$ and $q$ be Boolean-valued predicates
such that $p(a)=\neg q(a)$ for every $a$ in $l$. Then the listwise disjunction of $p$
over $l$ is the negation of the listwise conjunction of $q$ over $l$.
\end{theorem}

\begin{theorem}[Universal test and negated existential test]
\lean{BoolCircuit.all_eq_not_any}
\label{BoolCircuit.all_eq_not_any}
\leanok
\difficulty{4}
Let $l$ be a list of natural numbers and let $p,q$ be Boolean-valued functions on the
natural numbers. If $p(a)$ is the negation of $q(a)$ for every entry $a$ of $l$, then
every entry of $l$ satisfies $p$ exactly when no entry of $l$ satisfies $q$.
\end{theorem}

\begin{theorem}[Evaluation under duality]
\lean{BoolCircuit.DAGGate.eval_dual}
\label{BoolCircuit.DAGGate.eval_dual}
\leanok
\uses{BoolCircuit.DAGGate, BoolCircuit.DAGGate.dual, BoolCircuit.DAGGate.eval, BoolCircuit.GateKind.dual, BoolCircuit.all_eq_not_any, BoolCircuit.any_eq_not_all}
\difficulty{4}
Let a gate have all its input indices within the given Boolean value list, and suppose
that if the gate is a negation gate then it has exactly one input. Then evaluating its
dual gate on the list obtained by negating every value gives the negation of the
original gate's evaluation.
\end{theorem}

\begin{theorem}[Duality of gate execution]
\lean{BoolCircuit.runWith_dual}
\label{BoolCircuit.runWith_dual}
\leanok
\uses{BoolCircuit.DAGGate, BoolCircuit.DAGGate.dual, BoolCircuit.DAGGate.eval, BoolCircuit.DAGGate.eval_dual, BoolCircuit.runWith, BoolCircuit.runWith_cons}
\difficulty{5}
Let $G$ be a list of gates and $I$ a list of Boolean values. Assume that for each
position $i$ in $G$, every argument index of the gate at position $i$ is less than the
length of $I$ plus $i$, and that every NOT gate in $G$ has exactly one argument. Then
executing the gates obtained by interchanging AND and OR while retaining NOT and the
argument lists, starting from the entrywise negation of $I$, produces the entrywise
negation of the result of executing $G$ starting from $I$.
\end{theorem}

\begin{definition}[Dual circuit]
\lean{BoolCircuit.DAGCircuit.dual}
\label{BoolCircuit.DAGCircuit.dual}
\leanok
\uses{BoolCircuit.DAGCircuit, BoolCircuit.DAGGate.dual}
For a directed acyclic circuit on $n$ input vertices, its dual circuit has the same
output vertex and the same topological order of gates, with each gate replaced by its
dual. The circuit remains acyclic, and its output is a vertex among the $n$ inputs and
gates.
\end{definition}

\begin{theorem}[Evaluation under dualization]
\lean{BoolCircuit.DAGCircuit.eval_dual}
\label{BoolCircuit.DAGCircuit.eval_dual}
\leanok
\uses{BoolCircuit.DAGCircuit, BoolCircuit.DAGCircuit.dual, BoolCircuit.DAGCircuit.eval, BoolCircuit.DAGCircuit.values, BoolCircuit.runWith_dual}
\difficulty{4}
Let $C$ be a directed acyclic Boolean circuit on $n$ inputs in which every NOT gate has
exactly one argument. For every input assignment $x\in\{0,1\}^n$, the dual circuit
evaluated on the bitwise complement of $x$ outputs the complement of the output of $C$
on $x$.
\end{theorem}

\begin{theorem}[Dual preserves fan-in at most two]
\lean{BoolCircuit.DAGCircuit.IsFaninTwo.dual}
\label{BoolCircuit.DAGCircuit.IsFaninTwo.dual}
\leanok
\uses{BoolCircuit.DAGCircuit, BoolCircuit.DAGCircuit.IsFaninTwo, BoolCircuit.DAGCircuit.dual, BoolCircuit.DAGGate.dual, BoolCircuit.GateKind.dual}
\difficulty{4}
Let $C$ be a finite Boolean circuit whose gates have pairwise distinct arguments, whose
NOT gates have exactly one argument, and whose gates have at most two arguments. If the
dual circuit is formed by swapping AND and OR gates while leaving NOT gates and all
arguments unchanged, then it also has pairwise distinct gate arguments, exactly one
argument at each NOT gate, and at most two arguments at every gate.
\end{theorem}

\begin{definition}[Emitter states]
\lean{BoolCircuit.CktSatReduction.Emitter.PS}
\label{BoolCircuit.CktSatReduction.Emitter.PS}
\leanok
Let $PS$ be the finite set of seven emitter states $n,g,k_1,a,u,o$ and $k_2(b)$, where
$b$ ranges over the two Boolean values.
\end{definition}

\begin{definition}[Emitter state transition]
\lean{BoolCircuit.CktSatReduction.Emitter.psNext}
\label{BoolCircuit.CktSatReduction.Emitter.psNext}
\leanok
\uses{BoolCircuit.CktSatReduction.Emitter.PS}
The transition rule on emitter states and Boolean inputs is as follows: from $n$, input
$1$ leads to $n$ and input $0$ to $g$; from $g$, input $1$ leads to $k_1$ and input $0$
to $o$; from $k_1$, input $b$ leads to $k_2(b)$; from any $k_2(c)$, either input leads
to $a$; from $a$, input $1$ leads to $u$ and input $0$ to $g$; from $u$, input $1$ leads
to $u$ and input $0$ to $a$; and from $o$, either input leads to $o$.
\end{definition}

\begin{definition}[Conditional bit swap]
\lean{BoolCircuit.CktSatReduction.Emitter.swapBit}
\label{BoolCircuit.CktSatReduction.Emitter.swapBit}
\leanok
\uses{BoolCircuit.CktSatReduction.Emitter.PS}
For an emitter state $q$ and Boolean input $b$, output the complement of $b$ when
$q=k_2(0)$, and output $b$ in every other state.
\end{definition}

\begin{definition}[Bit transformation from the initial state]
\lean{BoolCircuit.CktSatReduction.swapKinds}
\label{BoolCircuit.CktSatReduction.swapKinds}
\leanok
\uses{BoolCircuit.CktSatReduction.Emitter.psNext, BoolCircuit.CktSatReduction.Emitter.swapBit, Complexity.transduce}
For a finite Boolean word $w$, process its bits in order using the emitter state
transition, starting in state $n$. At each step emit the input bit with its value
complemented exactly when the current state is $k_2(0)$; the resulting word is the
transformation of $w$.
\end{definition}

\begin{definition}[Output after entering the terminal state]
\lean{BoolCircuit.CktSatReduction.descTail}
\label{BoolCircuit.CktSatReduction.descTail}
\leanok
\uses{BoolCircuit.CktSatReduction.Emitter.psNext, Complexity.transduce}
For a finite Boolean word $w$, process its bits in order using the emitter state
transition, starting in state $n$, and emit a bit exactly when the current state is $o$,
emitting that input bit unchanged. The resulting word consists of the emitted bits.
\end{definition}

\begin{definition}[State dependent bit transformation]
\lean{BoolCircuit.CktSatReduction.T}
\label{BoolCircuit.CktSatReduction.T}
\leanok
\uses{BoolCircuit.CktSatReduction.Emitter.PS, BoolCircuit.CktSatReduction.Emitter.psNext, BoolCircuit.CktSatReduction.Emitter.swapBit, Complexity.transduce}
For an emitter state $q$ and a finite Boolean word $w$, process the bits of $w$ in order
using the emitter state transition, starting in state $q$. At each step emit the input
bit with its value complemented exactly when the current state is $k_2(0)$; the
resulting word is the transformation of $w$.
\end{definition}

\begin{definition}[State dependent terminal output]
\lean{BoolCircuit.CktSatReduction.S}
\label{BoolCircuit.CktSatReduction.S}
\leanok
\uses{BoolCircuit.CktSatReduction.Emitter.PS, BoolCircuit.CktSatReduction.Emitter.psNext, Complexity.transduce}
For an emitter state $q$ and a finite Boolean word $w$, process its bits in order using
the emitter state transition, starting in state $q$. Emit an input bit exactly when the
current state is $o$, and emit it unchanged; the resulting word consists of the emitted
bits.
\end{definition}

\begin{theorem}[Transduction of the empty string]
\lean{BoolCircuit.CktSatReduction.T_nil}
\label{BoolCircuit.CktSatReduction.T_nil}
\leanok
\uses{BoolCircuit.CktSatReduction.Emitter.PS, BoolCircuit.CktSatReduction.T}
\difficulty{1}
For every transducer state $q$, processing the empty bit string from $q$ returns the
empty bit string.
\end{theorem}

\begin{theorem}[Transduction of a leading bit]
\lean{BoolCircuit.CktSatReduction.T_cons}
\label{BoolCircuit.CktSatReduction.T_cons}
\leanok
\uses{BoolCircuit.CktSatReduction.Emitter.PS, BoolCircuit.CktSatReduction.Emitter.psNext, BoolCircuit.CktSatReduction.Emitter.swapBit, BoolCircuit.CktSatReduction.T}
\difficulty{1}
For every state $q$, bit $b$, and finite bit string $w$, processing $bw$ from state $q$
outputs the bit obtained by swapping $b$ according to $q$, followed by the result of
processing $w$ from the next state determined by $q$ and $b$:
$T_q(bw)=\bigl(	ext{swap}_q(b)\bigr)T_{	ext{next}(q,b)}(w)$.
\end{theorem}

\begin{theorem}[Output on the empty input]
\lean{BoolCircuit.CktSatReduction.S_nil}
\label{BoolCircuit.CktSatReduction.S_nil}
\leanok
\uses{BoolCircuit.CktSatReduction.Emitter.PS, BoolCircuit.CktSatReduction.S}
\difficulty{1}
Let the transducer have states $\{n,g,k_1,k_2(0),k_2(1),a,u,o\}$, transition function
given by the stated rules, and emit the input bit only when the current state is $o$,
emitting nothing otherwise. For every starting state $q$, its output on the empty bit
string is the empty string.
\end{theorem}

\begin{theorem}[Transducer output on an initial bit]
\lean{BoolCircuit.CktSatReduction.S_cons}
\label{BoolCircuit.CktSatReduction.S_cons}
\leanok
\uses{BoolCircuit.CktSatReduction.Emitter.PS, BoolCircuit.CktSatReduction.Emitter.psNext, BoolCircuit.CktSatReduction.S, Complexity.transduce}
\difficulty{3}
For the specified bit transducer, let $q$ be a state and let $b$ be a bit, followed by a
finite bit string $r$. The output on the input consisting of $b$ followed by $r$,
starting in state $q$, is the concatenation of the one-bit word $b$ if $q$ is the output
state (and the empty word otherwise) with the output on $r$ starting in the state
reached from $q$ upon reading $b$.
\end{theorem}

\begin{theorem}[Involution of the transducer]
\lean{BoolCircuit.CktSatReduction.T_T}
\label{BoolCircuit.CktSatReduction.T_T}
\leanok
\uses{BoolCircuit.CktSatReduction.Emitter.PS, BoolCircuit.CktSatReduction.Emitter.psNext, BoolCircuit.CktSatReduction.Emitter.swapBit, BoolCircuit.CktSatReduction.T, BoolCircuit.CktSatReduction.T_cons}
\difficulty{5}
For every state $q$ of the transducer and every finite bit string $w$, applying the
transducer twice, each time starting in state $q$, returns the original string:
$T_q(T_q(w))=w$.
\end{theorem}

\begin{theorem}[Involution of gate-kind swapping]
\lean{BoolCircuit.CktSatReduction.swapKinds_swapKinds}
\label{BoolCircuit.CktSatReduction.swapKinds_swapKinds}
\leanok
\uses{BoolCircuit.CktSatReduction.T_T, BoolCircuit.CktSatReduction.swapKinds}
\difficulty{1}
For every bit string $w$, applying twice the bit-string transformation that swaps AND
and OR gate kinds returns $w$.
\end{theorem}

\begin{theorem}[Preservation of leading ones]
\lean{BoolCircuit.CktSatReduction.T_ones}
\label{BoolCircuit.CktSatReduction.T_ones}
\leanok
\uses{BoolCircuit.CktSatReduction.Emitter.PS, BoolCircuit.CktSatReduction.Emitter.psNext, BoolCircuit.CktSatReduction.Emitter.swapBit, BoolCircuit.CktSatReduction.T, BoolCircuit.CktSatReduction.T_cons}
\difficulty{4}
Let $q$ be a transducer state such that reading a $1$ leaves the state unchanged and
outputs $1$. For every natural number $k$ and finite bit string $r$, processing $k$
consecutive $1$s followed by $r$ from state $q$ preserves those $k$ leading $1$s and
then processes $r$ from $q$.
\end{theorem}

\begin{theorem}[Ignoring repeated ones]
\lean{BoolCircuit.CktSatReduction.S_ones}
\label{BoolCircuit.CktSatReduction.S_ones}
\leanok
\uses{BoolCircuit.CktSatReduction.Emitter.PS, BoolCircuit.CktSatReduction.Emitter.psNext, BoolCircuit.CktSatReduction.S, BoolCircuit.CktSatReduction.S_cons}
\difficulty{4}
For the specified bit transducer, let $q$ be a state such that reading a one leaves the
state unchanged and $q$ is not the output state. For every natural number $k$ and finite
bit string $r$, starting in $q$ and reading $k$ consecutive ones followed by $r$
produces the same output as reading $r$ from $q$.
\end{theorem}

\begin{theorem}[Identity in the output state]
\lean{BoolCircuit.CktSatReduction.T_o}
\label{BoolCircuit.CktSatReduction.T_o}
\leanok
\uses{BoolCircuit.CktSatReduction.Emitter.psNext, BoolCircuit.CktSatReduction.Emitter.swapBit, BoolCircuit.CktSatReduction.T, BoolCircuit.CktSatReduction.T_cons}
\difficulty{4}
For every finite bit string $r$, processing $r$ from state $o$ returns $r$ unchanged.
\end{theorem}

\begin{theorem}[Output in the absorbing state]
\lean{BoolCircuit.CktSatReduction.S_o}
\label{BoolCircuit.CktSatReduction.S_o}
\leanok
\uses{BoolCircuit.CktSatReduction.Emitter.psNext, BoolCircuit.CktSatReduction.S, BoolCircuit.CktSatReduction.S_cons}
\difficulty{4}
For every finite bit string $r$, the transducer started in state $o$ outputs exactly
$r$.
\end{theorem}

\begin{theorem}[Transduction of encoded argument lists]
\lean{BoolCircuit.CktSatReduction.T_args}
\label{BoolCircuit.CktSatReduction.T_args}
\leanok
\uses{BoolCircuit.CktSatReduction.Emitter.psNext, BoolCircuit.CktSatReduction.Emitter.swapBit, BoolCircuit.CktSatReduction.T, BoolCircuit.CktSatReduction.T_cons, BoolCircuit.CktSatReduction.T_ones, BoolCircuit.encodeList, BoolCircuit.encodeNat}
\difficulty{4}
For every finite list $A$ of natural numbers and every finite bit string $r$, processing
the encoding of $A$ followed by $r$ from state $a$ leaves the encoded list unchanged and
processes the suffix from state $g$: $T_a(	ext{encode}(A)\rVert r)=	ext{encode}(A)\rVert
T_g(r)$.
\end{theorem}

\begin{theorem}[Argument encoding and transducer output]
\lean{BoolCircuit.CktSatReduction.S_args}
\label{BoolCircuit.CktSatReduction.S_args}
\leanok
\uses{BoolCircuit.CktSatReduction.Emitter.psNext, BoolCircuit.CktSatReduction.S, BoolCircuit.CktSatReduction.S_cons, BoolCircuit.CktSatReduction.S_ones, BoolCircuit.encodeList, BoolCircuit.encodeNat}
\difficulty{4}
For every finite list $\mathrm{args}$ of natural numbers and every finite bit string
$r$, the output of the transducer started in state $a$ on the concatenation of the
encoded list of unary encodings of $\mathrm{args}$ with $r$ equals its output when
started in state $g$ on $r$.
\end{theorem}

\begin{theorem}[Dualization of encoded gates]
\lean{BoolCircuit.CktSatReduction.T_gates}
\label{BoolCircuit.CktSatReduction.T_gates}
\leanok
\uses{BoolCircuit.CktSatReduction.Emitter.psNext, BoolCircuit.CktSatReduction.Emitter.swapBit, BoolCircuit.CktSatReduction.T, BoolCircuit.CktSatReduction.T_args, BoolCircuit.CktSatReduction.T_cons, BoolCircuit.CktSatReduction.T_o, BoolCircuit.DAGGate, BoolCircuit.DAGGate.dual, BoolCircuit.DAGGate.encode, BoolCircuit.GateKind.dual, BoolCircuit.GateKind.encode, BoolCircuit.encodeList}
\difficulty{4}
For every finite list $G$ of gates and every finite bit string $r$, processing the
encoding of $G$ followed by $r$ from state $g$ yields the encoding of the list obtained
by replacing each gate with its dual, followed by the unchanged suffix $r$.
\end{theorem}

\begin{theorem}[Output after an encoded gate list]
\lean{BoolCircuit.CktSatReduction.S_gates}
\label{BoolCircuit.CktSatReduction.S_gates}
\leanok
\uses{BoolCircuit.CktSatReduction.Emitter.psNext, BoolCircuit.CktSatReduction.S, BoolCircuit.CktSatReduction.S_args, BoolCircuit.CktSatReduction.S_cons, BoolCircuit.CktSatReduction.S_o, BoolCircuit.DAGGate, BoolCircuit.DAGGate.encode, BoolCircuit.GateKind.encode, BoolCircuit.encodeList}
\difficulty{4}
Let a Boolean gate be labelled AND, OR, or NOT and have a finite list of natural-number
input indices. Encode a natural number $k$ as $k$ copies of $1$ followed by $0$, encode
a list by placing $1$ before each encoded entry and ending with $0$, and encode each
gate by its two-bit label followed by the encoding of its input-index list. For every
finite list of such gates and every finite bit string $r$, the bit transducer that scans
these encodings, suppressing output until it reaches the suffix, outputs exactly $r$
when started in its gate-processing state and given the encoded gate list followed by
$r$.
\end{theorem}

\begin{theorem}[Encoding commutes with gate duality]
\lean{BoolCircuit.CktSatReduction.swapKinds_encode}
\label{BoolCircuit.CktSatReduction.swapKinds_encode}
\leanok
\uses{BoolCircuit.CktSatReduction.Emitter.psNext, BoolCircuit.CktSatReduction.Emitter.swapBit, BoolCircuit.CktSatReduction.T, BoolCircuit.CktSatReduction.T_cons, BoolCircuit.CktSatReduction.T_gates, BoolCircuit.CktSatReduction.T_ones, BoolCircuit.CktSatReduction.swapKinds, BoolCircuit.DAGCircuit, BoolCircuit.DAGCircuit.dual, BoolCircuit.DAGCircuit.encode, BoolCircuit.encodeNat}
\difficulty{3}
For every directed acyclic circuit $C$, applying the bit-string transformation that
swaps AND and OR gate kinds to the encoding of $C$ gives the encoding of the circuit
obtained by exchanging AND and OR at every gate while leaving NOT gates and all gate
arguments unchanged.
\end{theorem}

\begin{theorem}[Description tail of a circuit encoding]
\lean{BoolCircuit.CktSatReduction.descTail_encode}
\label{BoolCircuit.CktSatReduction.descTail_encode}
\leanok
\uses{BoolCircuit.CktSatReduction.Emitter.psNext, BoolCircuit.CktSatReduction.S, BoolCircuit.CktSatReduction.S_cons, BoolCircuit.CktSatReduction.S_gates, BoolCircuit.CktSatReduction.S_ones, BoolCircuit.CktSatReduction.descTail, BoolCircuit.DAGCircuit, BoolCircuit.DAGCircuit.encode, BoolCircuit.encodeNat}
\difficulty{3}
Let $C$ be a finite acyclic Boolean circuit with $n$ input vertices, whose gates are
listed in topological order with every gate reading only input or earlier vertices, and
whose output is a vertex of the circuit. The description-tail transducer applied to the
encoding of $C$ returns the unary encoding of the output vertex: a string of $C$'s
output-index many $1$s followed by $0$.
\end{theorem}

\begin{definition}[Validity of an encoded circuit]
\lean{BoolCircuit.CktSatReduction.isValid}
\label{BoolCircuit.CktSatReduction.isValid}
\leanok
\uses{BoolCircuit.DAGCircuit.IsFaninTwo, BoolCircuit.DAGCircuit.decode}
A finite Boolean word $w$ is valid if it decodes as some circuit $C$ and that circuit
has fan-in two; it is invalid if decoding fails or the decoded circuit does not have
fan-in two.
\end{definition}

\begin{theorem}[Characterization of valid circuit encodings]
\lean{BoolCircuit.CktSatReduction.isValid_iff}
\label{BoolCircuit.CktSatReduction.isValid_iff}
\leanok
\uses{BoolCircuit.CktSatReduction.isValid, BoolCircuit.DAGCircuit, BoolCircuit.DAGCircuit.IsFaninTwo, BoolCircuit.DAGCircuit.decode, BoolCircuit.DAGCircuit.decode_encode, BoolCircuit.DAGCircuit.encode, BoolCircuit.DAGCircuit.encode_of_decode}
\difficulty{4}
For every bit string $w$, its validity indicator is true if and only if there exist a
number $n$ of input vertices and a directed acyclic circuit on those inputs, with output
a circuit vertex, whose gates read only earlier vertices, whose argument lists have no
repeated entries, whose NOT gates have exactly one argument, and whose gates each have
at most two arguments, such that the circuit encodes to $w$.
\end{theorem}

\begin{theorem}[Input count bounded by encoding length]
\lean{BoolCircuit.CktSatReduction.le_length_encode}
\label{BoolCircuit.CktSatReduction.le_length_encode}
\leanok
\uses{BoolCircuit.DAGCircuit, BoolCircuit.DAGCircuit.encode, BoolCircuit.encodeNat}
\difficulty{1}
For every directed acyclic circuit $C$ with $n$ input vertices, the length of its
bit-string encoding is at least $n$.
\end{theorem}

\begin{theorem}[Validity by circuit evaluation]
\lean{BoolCircuit.CktSatReduction.isValid_eq_verdict}
\label{BoolCircuit.CktSatReduction.isValid_eq_verdict}
\leanok
\uses{BoolCircuit.CircuitEval.verdict, BoolCircuit.CircuitEval.verdict_pairEncode, BoolCircuit.CktSatReduction.isValid, BoolCircuit.CktSatReduction.isValid_iff, BoolCircuit.CktSatReduction.le_length_encode, BoolCircuit.CktSatReduction.swapKinds, BoolCircuit.CktSatReduction.swapKinds_encode, BoolCircuit.CktSatReduction.swapKinds_swapKinds, BoolCircuit.DAGCircuit.IsFaninTwo.dual, BoolCircuit.DAGCircuit.dual, BoolCircuit.DAGCircuit.encode, BoolCircuit.DAGCircuit.eval, BoolCircuit.DAGCircuit.eval_dual, Turing.pairEncode}
\difficulty{6}
For every bit string $w$, the validity indicator of $w$ equals the Boolean disjunction
of two verdicts: the verdict on the pairing of $w$ with a string of $|w|$ ones, and the
verdict on the pairing of the circuit-kind-swapped version of $w$ with a string of $|w|$
zeros.
\end{theorem}

\begin{theorem}[Polynomial-time circuit-encoding validity]
\lean{BoolCircuit.CktSatReduction.polyTimeComputable_isValid}
\label{BoolCircuit.CktSatReduction.polyTimeComputable_isValid}
\leanok
\uses{BoolCircuit.CircuitEval.evalTM, BoolCircuit.CircuitEval.evalTM_computes, BoolCircuit.CircuitEval.verdict, BoolCircuit.CktSatReduction.isValid, BoolCircuit.CktSatReduction.isValid_eq_verdict, BoolCircuit.CktSatReduction.swapKinds, Complexity.PolyTimeComputable, Complexity.PolyTimeComputable.append, Complexity.PolyTimeComputable.comp, Complexity.PolyTimeComputable.pairEncode, Complexity.polyTimeComputable_id, Complexity.polyTimeComputable_of_linear, Complexity.polyTimeComputable_transduce, Complexity.transduce, Turing.FinTM.computesFunInTime_ifEq, Turing.pairEncode}
\difficulty{5}
The function that maps each bit string $w$ to the one-bit string containing its validity
indicator is computable in polynomial time.
\end{theorem}
